\documentclass[11pt,a4paper]{article}

% --- arXiv-style preprint layout -------------------------------------------
\usepackage[margin=1in]{geometry}
\usepackage[T1]{fontenc}
\usepackage[utf8]{inputenc}
\usepackage{lmodern}
\usepackage{amsmath,amssymb}
\usepackage{booktabs}
\usepackage{graphicx}
\usepackage{xcolor}
\usepackage{url}
\usepackage{hyperref}
\usepackage{microtype}

\hypersetup{colorlinks=true, linkcolor=blue!60!black, citecolor=blue!60!black, urlcolor=blue!60!black}
\setlength{\parskip}{0.4em}
\setlength{\parindent}{0pt}

\newcommand{\hdr}[1]{\textbf{#1}}

\title{\bfseries Open Dream-RSI: An Open-Source Library for Recursive
Self-Improvement Around a Frozen LLM, with Replay-Gated Learned Policies and a
Curated Knowledge Base}

\author{%
Patryk Orwat\\
Independent Researcher\\
\texttt{patrykorwat@gmail.com}\\[2pt]
\small September 22, 2026
}
\date{}

\begin{document}
\maketitle

\begin{abstract}
We present \textbf{Open Dream-RSI}, an open-source (MIT) Python library that
implements the exploration layer of recursive self-improvement (RSI) around a
\emph{frozen} large language model: no weight access, no GPU, no training. The
library reproduces the discovery loop of Dream-RSI \cite{dreamrsi} and extends
it in three directions that the paper leaves implicit. First, an
\emph{always-on supervisor} that self-schedules improvement cycles under an
explicit API-call budget guard. Second, \emph{LLM-written exploration
policies}: the loop asks the model to write candidate expansion policies as
Python source, executes them in an isolated subprocess, scores them by
\emph{counterfactual replay rollout} on the recorded discovery tree, and
promotes them only on evidence. Third, a \emph{knowledge curator} that distils
verifier failures into short, schema-validated, deduplicated lesson records
retrieved into every future proposal---so iterations become not only cheaper
but measurably smarter. Fourth, \emph{thought-conditioned branching}: each
attempt records the one-line plan the model stated before its code, the
proposal prompt carries the task's tried-idea ledger, and expansion leaves a
branch as soon as its recorded idea repeats itself---semantic selection
replacing purely numeric score-greedy choice. We introduce two deterministic,
key-free benchmarks:
\texttt{bench}, which measures API-call savings (54\% at equal solve quality on
the built-in suite), and \texttt{bench-policy}, a \emph{decoy-trap} suite that
isolates exploration quality where solve-rate graphs saturate: score-greedy
exploration solves \textbf{0/120} tasks, $\varepsilon$-greedy 24\%, replay-gated
learned policies reach \textbf{72\%} (100\% by cycle~5) at \emph{fewer} total
calls than $\varepsilon$ spends on luck, the curated
knowledge base alone---with zero policy-generation calls---reproduces the same
72\%, and thought-guided branching reaches \textbf{100\% from the first cycle}
at 1.37 calls per solve even with exploration randomness disabled. All headline numbers are produced by a scripted solver and are therefore
claims about the \emph{machinery} (gates, sandboxes, promotion, retrieval),
not about any particular model; we state this limitation explicitly.
The package has zero runtime dependencies, $\sim$4k lines of library code,
1.1k lines of tests (81 passing), a live dashboard, and a single CLI.
\end{abstract}

\vspace{0.5em}
\hrule
\vspace{0.3em}
\small
\textbf{Code:} \url{https://github.com/patrykorwat/open-dream-rsi} ---
MIT license, Python $\geq$ 3.9, zero runtime dependencies, installable from
source with no build step.
\vspace{0.3em}
\hrule
\normalsize

\section{Introduction}

Recursive self-improvement systems that post-train model weights
\cite{openrsi} require GPU clusters, training infrastructure, and
weight access. The Dream-RSI line of work \cite{dreamrsi} shows a lighter
recipe: keep the model frozen and instead improve the \emph{process around
it}---an offline ``dreaming'' loop that optimises how the online agent
exploits accumulated discovery history. Since the paper's release
(2026-09-14) several independent implementations have appeared; Open
Dream-RSI is one of them, not the first (see \S\ref{sec:related}).

This library targets the gap between the paper's method sketch and a usable
tool: a researcher with an OpenAI-compatible endpoint should be able to
\texttt{pip}-install the loop, point it at their tasks and model, and get a
self-improving search that is (a)~budget-safe, (b)~safe to run code the model
writes, and (c)~capable of getting \emph{better}, not merely cheaper. We make
three contributions beyond the published method description:

\begin{enumerate}
  \item \hdr{Always-on supervisor with a hard budget guard.}
  \texttt{AutoRSIRuntime} self-schedules cycles, counts every API call
  (including offline dreaming), and refuses to overspend---the property any
  unattended deployment needs and no reference implementation ships.
  \item \hdr{LLM-written exploration policies with replay-gated promotion.}
  Per task category, the loop asks the frozen model for a policy program
  (\texttt{choose\_action(frontier, step)}), statically validates it, executes
  it in a scrubbed subprocess, scores it against the recorded discovery tree
  by counterfactual rollout, and promotes it only if it beats the incumbent on
  evidence (\S\ref{sec:policy}). Promotion, not trust, is the safety model.
  \item \hdr{Knowledge curator: lessons that make iterations smarter.}
  Recipes and policies capture \emph{what} worked; nothing captured \emph{why}
  attempts failed, so a trap re-pays its tuition every run. The curator
  distils verifier failures into short validated lesson records, retrieves
  them into future proposals, and prunes those that stop helping
  (\S\ref{sec:curator})---the skills/memory pattern of practical agent
  frameworks, applied to the loop itself.
\end{enumerate}

Two deterministic, key-free benchmarks back the claims (\S\ref{sec:bench}):
\texttt{bench} measures cost; \texttt{bench-policy} measures exploration on a
purpose-built decoy-trap suite, because on easy suites every strategy solves
everything and solve-rate graphs are blind to policy quality.

\section{System architecture}

The package is $\sim$3.5k lines of pure-Python with no runtime dependencies
(standard library only, including the LLM client) and 910 lines of tests.
Modules:

\begin{itemize}
  \item \texttt{core.tree} --- discovery tree: hypothesis, verifier score,
  action history per node.
  \item \texttt{core.agent} --- thin LLM-agent abstraction steered by the
  currently promoted policy.
  \item \texttt{core.simulator} / \texttt{core.dreamer} --- offline dreaming:
  replay and optimisation over recorded simulations without touching the
  online budget.
  \item \texttt{core.policygen} --- policy generation, validation, sandboxed
  execution, counterfactual rollout scoring, promotion gate
  (\S\ref{sec:policy}).
  \item \texttt{core.curator} --- lesson distillation, validation, retrieval,
  pruning (\S\ref{sec:curator}).
  \item \texttt{memory} --- persistent cross-run store: verified-solution
  recipes, promoted policies, \texttt{lessons.json}.
  \item \texttt{loop} --- \texttt{AutoRSIRuntime}: scheduling, budget guard,
  feedback wiring between dreaming and online expansion.
  \item \texttt{llm} --- OpenAI-compatible client
  (\texttt{POST \{base\_url\}/chat/completions}) via \texttt{urllib}; works
  against OpenAI, vLLM, Ollama, LM Studio, or any protocol-compatible server
  through \texttt{OPENAI\_BASE\_URL}.
  \item \texttt{tools} --- sandboxed code verification (same isolation
  boundary as the policy sandbox).
  \item \texttt{dashboard} --- dependency-free live web dashboard.
\end{itemize}

\hdr{Budget guard.} Every call---online proposals, offline dreaming, policy
generation, curation---increments a single counter checked before dispatch.
A cycle halts at its budget rather than overspending; the supervisor can then
sleep, re-plan, and resume from persisted memory. This is what makes the loop
deployable unattended: the cost per cycle is a declared quantity, not a hope.

\hdr{Sandbox boundary.} Candidate policy code and model-generated solutions
are never executed in-process. Execution happens in
\texttt{python~-I} subprocesses with a scrubbed environment (only
\texttt{PATH}), a hard timeout (5\,s per decision), and a temp-dir harness.
Invalid code is additionally rejected by AST-level static validation (shape:
\texttt{def choose\_action(frontier, step)}; no imports, no dunder access)
before any process spawns. API keys never cross the sandbox wall.

\section{LLM-written exploration policies}
\label{sec:policy}

For each task category the loop may request one LLM call producing a complete
policy program. The candidate then passes a pipeline in which every stage can
reject it:

\begin{enumerate}
  \item \hdr{Static validation.} AST parse; required entry point present;
  banned constructs (imports, \texttt{eval}, dunder access) reject outright.
  \item \hdr{Sandboxed execution} against recorded frontiers (above).
  \item \hdr{Counterfactual rollout scoring.} The candidate's decisions are
  replayed on the recorded discovery tree: at each simulated step the
  candidate picks a frontier node; the reward is the score the child actually
  achieved in the recorded ladder (back-propagated outcome for
  never-expanded nodes). One batched subprocess scores an entire rollout
  ($\sim$15$\times$ faster than per-step spawning). The objective is
  \begin{equation}
    \text{score}(\pi) \;=\; \frac{1}{T}\sum_{t=1}^{T} r_t
    \;+\; \beta_2 \cdot \frac{|\{\text{distinct picks}\}|}{T},
    \qquad \beta_2 = 0.2,
  \end{equation}
  applied \emph{identically} to the candidate and to an in-process greedy
  baseline. Mean-only scoring mis-ranks explorers below greedy on trap trees;
  the diversity term is what lets a genuinely structural explorer win the
  gate.
  \item \hdr{Promotion gate.} A candidate replaces the incumbent only with
  evidence: $\geq$80\% valid sandbox decisions (\texttt{MIN\_VALIDITY}=0.8)
  over a horizon of $\geq$8 recorded steps (\texttt{MIN\_HORIZON}), and a
  rollout score beating the greedy baseline under the shared objective.
  \item \hdr{Call economy.} A category that already solves with its promoted
  policy is not re-asked; regeneration requires 4 new replay steps of
  evidence (\texttt{POLICY\_REGROW\_STEPS}), and all generation calls count
  against the same budget guard.
\end{enumerate}

The promoted program then \emph{steers online expansion}: which frontier node
to expand next is decided by validated model-written code rather than a fixed
heuristic. This closes the loop the paper describes---the process around the
model improves---while every untrusted artefact (code and text alike) passes a
structural gate.

\section{The knowledge curator}
\label{sec:curator}

Recipes (verified solutions as warm starts) and policies capture \emph{what}
worked. Neither captures \emph{why} attempts fail, so the same trap re-buys
its tuition fee every run. The curator adds a third persistent artefact.

\hdr{Lesson records.} After failed verification, the loop invokes the model
once (one call, budget-counted) to distil the verifier evidence into lessons
carrying a short trigger phrase, a lesson text of $\leq$280 characters, and
$\leq$5 evidence snippets.

Lessons are \emph{pure text}---never executed---so their gate is structural
plus evidential rather than sandboxed: schema and caps (2--4 trigger words;
20--280 characters; $\leq$8 lessons per category; $\leq$4 lessons / 700
characters injected per prompt; $\leq$5 evidence snippets each), evidence
required at admission, and deduplication against the existing store.

\hdr{Retrieval feedback.} At proposal time the top lessons for the category
are retrieved (trigger-match, ranked by \texttt{wins}/\texttt{uses} counters
updated from subsequent verification outcomes); counters drive ranking and
eviction, so lessons that demonstrably stop helping are pruned automatically.
The knowledge base therefore cannot silently grow stale, and its worst
failure mode---injecting an unhelpful paragraph---is bounded by the caps.

\hdr{What is claimed and what is not.} The benchmark below demonstrates that
the \emph{wiring} works: a lesson, once admitted, causally changes what the
proposal prompt contains and the scripted world responds to it. The quality
of lesson \emph{writing} by a live model is an open question this library
deliberately does not claim to have settled (\S\ref{sec:limits}).

\section{Thought-conditioned branching}
\label{sec:thoughts}

Policies (section~\ref{sec:policy}) select branches by recorded scalars only:
the frontier observation carries scores, visit counts and error strings, but
nothing about \emph{which idea} a node stands for. A loop can therefore
re-expand a branch forever without knowing it is re-polishing the same dead
plan, and a proposal cannot see which ideas the branch already exhausted.
Thought-conditioned branching (the runtime default,
\texttt{enable\_thoughts=True}, ablatable with \texttt{--no-thoughts}) closes
this gap with three cheap, purely textual additions.

\hdr{Per-node idea label.} When enabled, the candidate system prompt asks the
model for exactly one \texttt{PLAN:} line before the code block; the line is
extracted tolerantly (bullets, dropped colons and case are forgiven) and
stored as \texttt{TreeNode.thought}, truncated at 240 characters. Like
lessons, thoughts are pure text---never executed, never fed to a sandbox---so
their admission requires no gate beyond length, and they are persisted with
the archived tree, so replay worlds and later cycles see the same labels.
The frontier entry and the policy-generation summary gain a \texttt{thought}
field, letting LLM-written policies reason over idea families, not just scores.

\hdr{Tried-idea ledger.} Every proposal prompt additionally carries the
task's deduplicated idea ledger---best recorded outcome per distinct idea,
$\leq$8 lines---rendered above the policy hints: which plans have already
been attempted in this task, and what the verifier said about each. The
model is instructed not to repeat a dead idea verbatim and to build on
structurally promising ones. The ledger is deliberately task-wide rather
than branch-local: whether an idea is dead does not depend on which node
happens to be expanded next.

\hdr{Semantic escape from idea stagnation.} At expansion time, the loop
applies the rule shown below after the promoted policy program declines to
act and before the $\varepsilon$/greedy cascade decides:

\begin{quote}\small\begin{minipage}{0.94\linewidth}
\textbf{Stagnation steering.} Let $b$ be the score-greedy frontier pick with
idea $t$. If $t$ is empty, or no recorded child of $b$ carries $t$ again
(the branch is converging, not repeating), do nothing---normal cascade.
Otherwise $b$ is polishing a dead plan: choose the frontier node
maximising $(\mathit{outcome}, -\mathit{children})$ among (i) nodes labelled
with a \emph{different}, non-empty idea, or, if none exist, (ii) nodes with
no idea label at all (typically the seed, whose ladder may still hold
untried ideas). No randomness, no API call, no verifier score improving:
the trigger is recorded semantic evidence about the \emph{idea}, and the
selection prefers the least-expanded node with the best known continuation.
\end{minipage}\end{quote}

The rule strictly refines the existing machinery: when nothing is stagnant
it is inert, and when it fires it can only redirect expansion to a node the
policy cascade was already allowed to choose---so the worst outcome of a
mislabelled thought is neutral, matching the fallback philosophy that
protects the section-3 gate.

\hdr{Causal attribution.} The decoy-trap benchmark's fifth arm
(\S\ref{sec:bench}) isolates the contribution to the \emph{expansion-time
steering}, not to the prompt text: with $\varepsilon$ forced to zero and
both policy generation and curation disabled the arm escapes every trap
from the first cycle (120/120, 164 calls, 1.37 calls per solve), while the
\texttt{thought\_steering=False} ablation (PLAN lines still recorded and the
ledger still rendered into proposals, semantic pick disabled) collapses to
the baselines---0/120 under greedy expansion, identical to
$\varepsilon$-greedy's 29/120 with $\varepsilon$ on. The scripted solver is
deliberately no more generous to the thought arm than to the lesson arm:
the ledger may redirect a repeatedly-polished dead idea to the alternative
idea, never directly to the fix, and tests assert that every solved task's
fix node grows from expanding the promising-idea node---the same ladder
every other arm must climb. Selection therefore stops being purely numeric
on evidence: the loop knows \emph{which idea} a branch stands for, and
abandons it when the record says the idea is dead.

\section{Benchmarks}
\label{sec:bench}

Both benchmarks are deterministic and key-free: a scripted solver makes the
model arm a constant, so every delta is attributable to the library's
machinery. They are reproducible from a fresh clone in seconds:

\begin{verbatim}
python -m open_dream_rsi bench --cycles 10 --markdown
python -m open_dream_rsi bench-policy --cycles 8 --format md
\end{verbatim}

\subsection{Cost: the dreaming loop against a cold baseline}

\texttt{bench} runs the built-in suite (10 cycles $\times$ 5 tasks) twice:
a cold baseline (fresh memory each cycle) and the full loop (dreaming +
persistent recipes and policies), same scripted model. Measured at the
defaults:

\begin{table}[h]
\centering\small
\begin{tabular}{lrrr}
\toprule
arm & solves & API calls & calls / task$\cdot$cycle \\
\midrule
cold baseline (fresh memory each cycle) & 50 & 150 & 3.00 \\
dream loop (policies + recipes + dreaming) & 50 & 69 & 1.38 \\
\bottomrule
\end{tabular}
\caption{\texttt{bench}: 54.0\% API-call saving at equal solve quality
(3{,}000 dreaming iterations run offline; wall-clock 1.6\,s). The curator is
disabled in this benchmark to keep the cost measurement pure.}
\end{table}

\subsection{Exploration: the decoy-trap suite}

Solve-rate graphs cannot show that a learned policy helps, because easy
suites saturate every arm. \texttt{bench-policy} constructs 5 tasks whose fix
hides behind a \emph{low-scoring} branch (passes 1/3 tests) past a plausible
\emph{decoy} that passes 2/3 forever; family scores are fixed by construction
and asserted by tests. Score-greedy expansion locks onto the decoy
permanently; escaping requires either structural exploration or memory.
Four arms share one scripted solver and budget and differ only in machinery
(3 runs $\times$ 8 cycles $\times$ 5 tasks, 24 calls/cycle):

\begin{table}[h]
\centering\small
\begin{tabular}{lrrrrr}
\toprule
arm & solves & rate & calls & calls/solve & gen.\ calls \\
\midrule
greedy & 0/120 & 0\% & 480 & $\infty$ & 0 \\
$\varepsilon$-greedy ($\varepsilon{=}0.3$) & 29/120 & 24\% & 405 & 13.97 & 0 \\
replay-gated policies & \textbf{86/120} & \textbf{72\%} & \textbf{274} & \textbf{3.19} & 33 policy \\
curated knowledge base & \textbf{87/120} & \textbf{72\%} & 283 & 3.25 & 40 curator \\
thought-guided branching & \textbf{120/120} & \textbf{100\%} & \textbf{164} & \textbf{1.37} & 0 \\
\bottomrule
\end{tabular}
\caption{\texttt{bench-policy}: the knowledge arm is
$\varepsilon$-greedy \emph{with zero policy calls}---every escape above the
$\varepsilon$ baseline is attributable to remembered lessons. Policies and
knowledge reach the same plateau (100\% by cycle~5 and~8 respectively);
neither is ranked over the other, and 87 vs.\ 86 is noise. The thought arm
additionally passes the \texttt{thought\_steering=False} ablation: with the
semantic expansion pick disabled it falls back to 0/120 (greedy) and 29/120
($\varepsilon$), pinning its win to expansion-time steering (\S\ref{sec:thoughts}).}
\end{table}

\begin{figure}[h]
\centering
\includegraphics[width=\linewidth]{odr_policy_bench.png}
\caption{Per-cycle solve rate on the decoy-trap suite (mean over 3 seeds).
Greedy never escapes the decoy; $\varepsilon$-greedy escapes by luck;
replay-gated policies and the curated knowledge base saturate at 100\% only
after several cycles; thought-guided branching---the shipped library
default---escapes every trap from the first cycle at the lowest total cost.}
\end{figure}

\hdr{Seed robustness.} The shipped seeds (7,\,11,\,23) are not cherry-picked:
over five independent seed sets the two knowledge/policy arms land within
78--88 and 78--87 solves of 120 respectively---statistically
indistinguishable from each other---while $\varepsilon$-greedy stays at
21--32. The dominance over $\varepsilon$ (roughly 3$\times$ the solves at
$\sim$70\% of its calls, 3.2 vs.\ 14 calls per solve) holds on every set.
The thought arm reaches 120/120 on both three-run seed sets and 159/160 and
200/200 on the four- and five-run sets---its per-set rate never drops
below 99\% and stays above the other machinery arms on every set.

\subsection{Test suite}

66 tests pass (910 LOC), including the benchmark's behavioural contracts:
trap-family scores equal their constructed values; the greedy arm must solve
nothing; the gated-policy arm must dominate $\varepsilon$; curator schema,
caps, dedup, and prune logic are pinned by unit tests.

\section{Related work and positioning}
\label{sec:related}

\hdr{Discovery loops around frozen LLMs.} FunSearch \cite{funsearch} pairs a
pretrained LLM with a systematic evaluator and evolves \emph{programs} rather
than solutions, yielding new results in extremal combinatorics and bin
packing; AlphaEvolve \cite{alphaevolve} generalises this to an evolutionary
coding agent over a program database grounded by automatic evaluation. The
G\"odel-machine ideal \cite{godel}---provably beneficial self-modification---is
relaxed by the Darwin G\"odel Machine \cite{dgm}, which empirically validates
each self-modification and grows an archive of agents. Dream-RSI
\cite{dreamrsi} (Zheng et al., Google/DeepMind/UMD, arXiv:2609.14858) is the
method this library implements; the official repository is
\url{https://github.com/zhengkid/Dream-RSI}. Our contribution is at the
library/infrastructure level within this family: the mechanisms above are
proven at scale by their authors; we ship the same loop pattern as a
dependency-free, budget-guarded, OpenAI-compatible tool and add the curator
and the exploration-isolating benchmark.

\hdr{Memory and self-generated policies in agents.} Reflexion
\cite{reflexion} shows that verbal self-feedback stored in an episodic memory
materially improves subsequent attempts---the same principle our knowledge
curator operationalises persistently across runs; Voyager \cite{voyager}
builds a persistent skill library of verified programs, the closest ancestor
of our replay-gated policy promotion.

\hdr{Independent implementations.} Since the paper's release several
independent implementations have appeared: a model-agnostic adapter SDK with
strict replay and evidence-based promotion gates \cite{drsdk}, a local-Qwen
CUDA-kernel exploration on DGX Spark \cite{dspark}, and agent/skill variants
\cite{hermesdrsi}. Against peers, this library differentiates on: the
always-on budget-guarded supervisor; LLM-\emph{written} policy code (as
opposed to a fixed policy menu) with counterfactual-replay promotion; a
dedicated exploration benchmark (\texttt{bench-policy}); a curated lesson
knowledge base with retrieval feedback; persistent cross-run memory; a live
dashboard; zero dependencies; MIT license.

\hdr{A different layer.} OpenRSI/OpenMLE/Frontis-MA1 \cite{openrsi,
mlebench} post-train weights (SFT+RL, a 35B meta-evolution agent on
MLE-Bench); the two layers are complementary---their trained improver could
be the frozen model behind our client, our dreaming loop could sit on top of
their search.

\section{Limitations}
\label{sec:limits}

We state plainly what the numbers do and do not show.
\begin{enumerate}
  \item \hdr{Scripted worlds, not model quality.} Every headline number comes
  from deterministic benchmarks with a scripted solver; the trap scores are
  fixed by construction, the scripted solver treats a stored lesson text as a
  genuine insight when it is visible in the prompt, and the ``policy''
  returned by the scripted model is a canned program. The benchmarks
  therefore validate the \emph{machinery}---sandboxes, gates, promotion,
  retrieval, budget---and prove the arms differ \emph{causally} through
  memory, not that any live model writes good lessons or policies. Live-model
  evaluation is the required next experiment.
  \item \hdr{Benchmark self-design risk.} The decoy-trap suite was designed by
  the same author as the loop it measures. Mitigations: family scores fixed
  and asserted by construction; arm contracts enforced by tests (greedy must
  fail, arms must differ causally); determinism lets anyone re-measure with
  one command; five seed sets reported.
  \item \hdr{Single-node sandbox.} \texttt{python -I} with scrubbed environment
  and timeout is a real boundary against casual escape but is not a container;
  deployments with hostile task sources should add OS-level isolation.
  \item \hdr{Scope.} No weight updates, no multi-agent dynamics, no adversarial
  red-teaming of promoted policies; the knowledge base is per-category text,
  not a structured world model.
\end{enumerate}

\section{Conclusion}

Open Dream-RSI turns the Dream-RSI method into a deployable, budget-safe,
dependency-free library and adds the two mechanisms a running loop actually
needs to compound---evidence-gated learned policies and a curated,
self-pruning knowledge base---each backed by a deterministic benchmark that
isolates exactly the capability it claims. At library scale: 54\% fewer API
calls at equal solve quality; on a suite where greedy exploration solves
nothing, replay-gated policies and remembered knowledge both reach 72\%
cumulative (100\% by cycle~5--8) at $\sim$4$\times$ better cost-per-solve
than random escape. All artefacts pass structural gates before touching the
loop; the gate, not trust, decides.


\begin{thebibliography}{99}
\bibitem{dreamrsi} Zheng et al., \emph{Dream-RSI: Recursive
Self-Improvement through Evolving Worlds}, arXiv:2609.14858, 2026.
\bibitem{funsearch} B.~Romera-Paredes, M.~Barekatain, A.~Novikov, M.~Balog,
M.~P.~Kumar, E.~Dupont, F.~J.~R.~Ruiz, J.~S.~Ellenberg, P.~Wang, O.~Fawzi,
P.~Kohli, A.~Fawzi, \emph{Mathematical discoveries from program search with
large language models}, Nature 625(7995), 2024.
DOI:~10.1038/s41586-023-06924-6.
\bibitem{alphaevolve} A.~Novikov, N.~V\~u, M.~Eisenberger, E.~Dupont,
P.-S.~Huang, A.~Z.~Wagner, S.~Shirobokov, B.~Kozlovskii, et~al.,
\emph{AlphaEvolve: A coding agent for scientific and algorithmic discovery},
arXiv:2506.13131, 2025.
\bibitem{dgm} J.~Zhang, S.~Hu, C.~Lu, R.~Lange, J.~Clune, \emph{Darwin
G\"odel Machine: Open-Ended Evolution of Self-Improving Agents},
arXiv:2505.22954, 2025.
\bibitem{godel} J.~Schmidhuber, \emph{G\"odel Machines: Fully
Self-Referential Optimal Universal Self-Improvers}, in B.~Goertzel and
C.~Pennachin (eds.), \emph{Artificial General Intelligence}, Springer, 2006.
\bibitem{reflexion} N.~Shinn, F.~Cassano, A.~Gopinath, K.~Narasimhan,
S.~Yao, \emph{Reflexion: Language Agents with Verbal Reinforcement
Learning}, Advances in Neural Information Processing Systems 36, 2023.
arXiv:2303.11366.
\bibitem{voyager} G.~Wang, Y.~Xie, Y.~Jiang, A.~Mandlekar, C.~Xiao, Y.~Zhu,
L.~Fan, A.~Anandkumar, \emph{Voyager: An Open-Ended Embodied Agent with
Large Language Models}, arXiv:2305.16291, 2023.
\bibitem{openrsi} FrontisAI, \emph{OpenRSI / OpenMLE / Frontis-MA1},
\url{https://github.com/FrontisAI/OpenRSI}.
\bibitem{mlebench} J.~S.~Chan, N.~Chowdhury, et~al., \emph{MLE-bench:
Evaluating Machine Learning Agents on Machine Learning Engineering},
arXiv:2410.07095, 2024.
\bibitem{drsdk} TheAstrayDev, \emph{dream-rsi-sdk},
\url{https://github.com/TheAstrayDev/dream-rsi-sdk}.
\bibitem{dspark} robinber, \emph{dream-rsi-spark},
\url{https://github.com/robinber/dream-rsi-spark}.
\bibitem{hermesdrsi} lesterppo, \emph{hermes-dream-rsi},
\url{https://github.com/lesterppo/hermes-dream-rsi}; Harkit2004,
\dots \emph{dream-rsi-skill}; mailbobg, \emph{Pi-RSI}.
\end{thebibliography}

\end{document}
